\section{Mono3DV}
\begin{figure*}[h!]
    \centering
    \includesvg[width=\textwidth]{images/mono3dv/overall.svg}
    \caption{\textbf{The overall of our proposed framework Mono3DV.} The architecture initially extracts features from a single-view image using an image backbone, which are then fed into a Transformer encoder. The subsequent decoder utilizes both standard learnable queries and supplementary noisy queries generated by a Variational Query Generator. Finally, the loss is determined by subjecting the predictions derived from the learnable queries to 3D-Aware Bipartite Matching.}
    \label{fig:mono3dv_overall}
\end{figure*}

\subsection{Overview}
\label{mono3dv:overview}
\Cref{fig:mono3dv_overall} illustrates the architecture of our proposed framework. Given a single-view image, we first employ ResNet50~\cite{ResNet} to extract high-level feature maps, which are subsequently input to a transformer encoder. Following the baseline MonoDETR\cite{MonoDETR}, the decoder utilizes $G$ groups of learnable queries $Q_L = \{q_{L_i}\}_{i=1}^G$, where $q_{L_i}\in\mathbb{R}^{N \times D}$, with $N$ representing the number of queries in each group and $D$ denoting the hidden dim. For each of these $G$ groups of learnable queries, we introduce $C$ corresponding groups of noisy queries $Q_N = \{\{q_{N_{ij}}\}_{j=1}^C\}_{i=1}^G$, where for the $i$-th group of learnable queries, we have $C$ associated groups of noisy queries, and $q_{N_{ij}}\in\mathbb{R}^{K \times D}$, with $K$ signifying the number of objects present in the input image. These noisy queries are generated by a Variational Query Generator to overcome the gradient vanishing issue associated with conventional denoising. During training, the predictions derived from the learnable queries are passed to 3D-Aware Bipartite Matching, which addresses the mismatch issues often encountered with naive 2D-only approaches.

\subsection{3D-Aware Bipartite Matching}
\textbf{3D Matching.} To correctly associate each query with its corresponding ground-truth object, we introduce a novel matching cost function, which is formally defined as:
\begin{equation}
    C_{match}=C_{2D} + \Gamma(t) C_{3D}
    \label{eq:matching_eq}
\end{equation}
The total matching cost is a weighted combination of 2D and 3D prediction costs. Specifically, $C_{2D}$ aggregates costs related to the object category, 2D bounding box size, and the projected 3D center. We follow the baseline MonoDETR~\cite{MonoDETR} and constructed the 2D cost as:
\begin{equation}
    \begin{aligned}
        C_{2D} & =\lambda_{cls}C_{cls} + \lambda_{proj}C_{xy3D} \\
        & \hspace{2em} + \lambda_{lrtb}C_{lrtb} + \lambda_{GIoU}C_{GIoU}
    \end{aligned}
\end{equation}
with $\lambda_{cls}=2, \lambda_{proj}=10, \lambda_{lrtb} = 5, \lambda_{GIoU} = 2$.

Conversely, $C_{3D}$ encompasses costs for depth, 3D size, and orientation angle. The 3D cost was formulated by defining a cost for each of the three 3D attributes, specifically size, angle, and depth, the detailed formulation of which is defined as follows:
\begin{equation}
    C_{3D} = C_{size3D} + C_{orien} + C_{depth}
\end{equation}

$\Gamma(t) \in [0, 1]$ serves as a scheduler weight. This weight is designed to be low during the initial training phase when 3D predictions are inherently unstable, and it gradually increases as the 3D estimation becomes more reliable.

\noindent\textbf{3D Weight Scheduler.} During initial training, the instability of the predicted 3D attributes makes their immediate integration into the Bipartite Matching problematic. We address this unreliability by introducing a step scheduler. This scheduler controls and delays the integration of the 3D attributes until prediction stability improves, which is formally defined as:
\begin{equation}
      \Gamma(t)=
      \left\{
      \begin{array}{l}
           0, \texttt{ if } t < T \\
           \epsilon, \texttt{ otherwise}
      \end{array}
      \right.
      \label{eq:3d_weight_scheduler}
\end{equation}
where $t$ denotes the current training epoch, $\epsilon$ is the weight threshold and $T$ is the trigger epoch.

\subsection{3D-DeNoising}
\label{mono3dv:3d_denoising}
The use of a scheduler weight, $\Gamma(t)$, when integrating 3D attributes into the Bipartite Matching process is designed to mitigate the mismatch problem in 2D-only approach. However, the potential benefit is severely compromised by the instability introduced by the ill-posed nature of 3D estimation from monocular images. This fundamental limitation restricts the method's impact, leading to performance that is only marginally better than the 2D-only approach. To overcome this limitation and decisively stabilize the 3D integration, we propose 3D-Denoising. By attaching the 3D ground truth to the noisy query generator, we inject strong, reliable 3D supervision directly into the training phase. This supervision acts as a robust anchor, guiding the network to learn a more stable and accurate 3D representation despite the noisy inputs, thereby stabilizing the benefits of the 3D attribute integration into the Bipartite Matching.

\noindent\textbf{3D Noisy Query.} In 2D object detection, DN-DETR~\cite{DN-DETR} utilizes the bounding box as the reference anchor and category of object to generate the noise query. We first reformulate how to generate a noisy query to enhance the effect of integrating 3D-attribute into Bipartite Matching. The objective of monocular 3D object detection consist of: category $c$, projected center $\left(x_c,y_c\right)$, 2D bounding box $l,r,t,b$, 3D dimension $l_{3D},w_{3D},h_{3D}$, orientation $\theta$ and central depth $d$.

To begin the denoising process, we first generate noisy bounding box information derived from the original ground truth. We introduce parameters $\lambda_C$ and $\lambda_D$ to manage the extent of applied noise.

\textbf{Projected 3D Center} $\left(x_c,y_c\right)$. Inspired by DN-DETR~\cite{DN-DETR}, we perform center shifting on the projected 3D center so that it still lies inside the original 2D bounding box by following the formula:
\begin{equation}
     \begin{aligned}
         \widetilde{x}_c &= x_c + \lambda_C \mathcal{U}(-1, 1) \cdot \frac{r+l}{2} \\
         \widetilde{y}_c &= y_c + \lambda_C \mathcal{U}(-1, 1) \cdot \frac{t+b}{2}
     \end{aligned}
\end{equation}

\textbf{2D Size} $\left(l,r,t,b\right)$. We also perform box scaling on the 2D bounding box, by manually adjusting $l,r,t,b$ through following formulation:
\begin{equation}
    \begin{aligned}
        \widetilde{l} &= \texttt{clip}\left(l + \lambda_C \mathcal{U}(-1, 1) \cdot l, 0, 1\right) \\
        \widetilde{r} &= \texttt{clip}\left(l + \lambda_C \mathcal{U}(-1, 1) \cdot r, 0, 1\right) \\
        \widetilde{t} &= \texttt{clip}\left(l + \lambda_C \mathcal{U}(-1, 1) \cdot t, 0, 1\right) \\
        \widetilde{b} &= \texttt{clip}\left(l + \lambda_C \mathcal{U}(-1, 1) \cdot b, 0, 1\right) \\
    \end{aligned}
\end{equation}
with \texttt{clip()} operation to limit the noisy value in range $\left[0,1\right]$.

\textbf{3D Size} $\left(l_{3D},h_{3D},w_{3D}\right)$. Inspired by established 2D bounding box scaling techniques, we extend this concept to 3D dimensions, employing the following formulation:
\begin{equation}
    \begin{aligned}
        \widetilde{l}_{3D} &= l_{3D} + \lambda_C \mathcal{U}(-1, 1) \cdot l_{3D} \\
        \widetilde{h}_{3D} &= h_{3D} + \lambda_C \mathcal{U}(-1, 1) \cdot h_{3D} \\
        \widetilde{w}_{3D} &= w_{3D} + \lambda_C \mathcal{U}(-1, 1) \cdot w_{3D}
    \end{aligned}
\end{equation}

\textbf{Object Depth} $d$. Noisy depth is generated by randomly adjusting the center depth along the object's length, with the condition that the resulting noisy depth value stays within the 3D bounding box, formulated as:
\begin{equation}
    \begin{aligned}
        \widetilde{d} &= d + \lambda_C \mathcal{U}
        (-1, 1) \cdot \frac{l_{3D}}{2}
    \end{aligned}
\end{equation}

\textbf{Category} $c$. By adopting label flipping~\cite{DN-DETR} for label noising with flipping rate $\lambda_D$, we aim to improve the model's ability to learn the relationship between object labels and their corresponding noisy bounding box predictions.

\textbf{Orientation} $\theta$. The baseline MonoDETR~\cite{MonoDETR} represents orientation by dividing it into multiple discrete bins, denoted as $\theta_{bin}$, and a continuous residual, denoted as $\theta_r$. To introduce noise into the orientation, we retain the residual $\theta_r$ and perform bin flipping with flipping rate $\lambda_D$. This involves randomly assigning the ground truth orientation bin to a different bin. Subsequently, only this noisy orientation bin is fed into the embedding layer, while the corresponding residual $\theta_r$ is discarded. This process ensures that information about the precise orientation is not entirely lost, as the residual component is preserved.

\textbf{Noisy Boxes Embedding.} Given the noisy continuous value $\left(\widetilde{l}_{3D}, \widetilde{h}_{3D}, \widetilde{w}_{3D}, \widetilde{d}\right)$, we apply sinusoidal positional encoding to map it into a high-dimensional vector. For the noisy discrete category $\widetilde{c}$ and orientation bin $\widetilde{\theta}_{bin}$, we utilize learnable embeddings. These encoded continuous features and discrete embeddings are then concatenated and subsequently fed into a three-layer Multi-Layer Perceptron (MLP) to obtain the noisy query vectors.

\textbf{3D Noisy Query Generator.} As described in Section 3.1, we introduce $Q_N = \{\{q_{N_{ij}}\}_{j=1}^C\}_{i=1}^G$ as noisy queries for the denoising process, where $q_{N_{ij}}\in \mathbb{R}^{K\times D}$ with $K$ being the number of objects in the input image. Prior to the addition of noise and feeding into the embedding layer, the ground truth object set is first repeated $C \times G$ times to obtain the exact number of queries needed. After generating noise boxes, we denote a 6D anchor box $\left(x_c,y_c,l,r,t,b\right)$ as the initial reference. Then we map the 3D information $\left(c,l_{3D},w_{3D},h_{3D},\theta,d\right)$ to a hidden space using an embedding layer, yielding the 3D query. The anchor box and the query are fed into the decoder to be reconstructed.

\noindent\textbf{Mask Separated Self-Attention.} To improve detection accuracy, the baseline MonoDETR model incorporates a group-wise one-to-many assignment strategy~\cite{GroupDETR}, employing $G$ distinct groups of learnable queries $Q_L = \{q_{L_i}\}_{i=1}^G$. This approach utilizes Separate Self-Attention~\cite{GroupDETR} to ensure that queries belonging to different groups do not interact, thereby maintaining group independence. To leverage the benefits of both the denoising training paradigm and the group-wise assignment methodology, we introduce the Mask Separated Self-Attention mechanism.

This mechanism employs a predefined attention $M \in \{0,1\}^{S \times S}$ inspired by DN-DETR~\cite{DN-DETR}, where $S = K \cdot C + N$, to regulate the interactions only among noisy queries, as well as between noisy queries and learnable queries. Adhering to the rational constraint of preventing information leakage as in DN-DETR~\cite{DN-DETR}, we devise an attention mask that satisfies two conditions: (1) learnable queries must not aggregate noisy queries to ensure consistency between training and inference, given that noisy queries are discarded during inference; (2) noisy queries within each group must not attend to those of other groups.
The detailed design of $M$ can be expressed by the following formulation:
\begin{equation}
    M_{i,j} = 
     \left\{
     \begin{array}{l}
          1, \texttt{ if } j < K\times C \texttt{ and } \lfloor\frac{i}{C}\rfloor \neq \lfloor\frac{j}{C}\rfloor \\
          1, \texttt{ if } j < K\times C \texttt{ and } i \geq K\times C\\
          0, \texttt{ otherwise}
     \end{array}
     \right.
\end{equation}
where $M_{i,j}=1$ means the i-th query cannot access the j-th query and $M_{i,j}=0$ otherwise.

Specifically, from the sets of learnable queries $Q_L$ and noisy queries $Q_N$, we construct a new set of combined queries $Q = \{q_i\}_{i=1}^G$, where each $q_i \in \mathbb{R}^{S \times D}$ is constructed by concatenating the learnable query $q_{L_i}$ with its corresponding noisy queries $\{q_{N_{ij}}\}_{j=1}^C$:
\begin{equation}
    q_i = \text{concat}( q_{N_{i1}}, q_{N_{i2}}, ..., q_{N_{iC}},q_{L_i})
    \label{eq:q_i}
\end{equation}
This newly formed set of queries is then processed by Separated Self-Attention with the predefined mask $M$. Through this proposed Mask Separated Self-Attention, we achieve precise control over the interaction patterns between queries of different types and within each group of queries. The implementation details are shown in Alg.~\ref{alg:mask_separated_self_attention}.

\begin{algorithm}[h]
\caption{Pseudocode of Mask Separated Self-Attention}
\label{alg:mask_separated_self_attention}
\definecolor{codeblue}{rgb}{0.25,0.5,0.5}
\definecolor{codekw}{rgb}{0.85, 0.18, 0.50}
\lstset{
  backgroundcolor=\color{white},
  basicstyle=\fontsize{7.5pt}{7.5pt}\ttfamily\selectfont,
  columns=fullflexible,
  breaklines=true,
  captionpos=b,
  commentstyle=\fontsize{7.5pt}{7.5pt}\color{codeblue},
  keywordstyle=\fontsize{7.5pt}{7.5pt}\color{codekw},
}
\begin{lstlisting}[language=python]
# SA: Self-Attention in the decoder layer
# Q_L: learnable queries, with size (GxN, B, D)
# Q_N: noisy queries, with size (G,CxK,B,D)
# M: attention mask, with size (S,S)
# B: batch size

# Mask Separated Self-Attention
if training:
    # Construct a new set of queries Q
    Q_L = Q_L.split(N, dim = 0) # a list of G tensors with shape (N,B,D)
    Q_L = cat(Q_L,dim=1) # (N,BxG,D)
    Q_N = Q_N.split(1, dim = 0).squeeze(0) 
    # a list of G tensors with shape (CxK,B,D)
    Q_N = cat(Q_N,dim=1) # (CxK,BxG,D)
    Q = cat((Q_N,Q_L),dim=0) # (CxK+N,BxG,D) or (S,BxG,D)

    # mask self-attention
    out = SA(Q,M) # (S,BxG,D)

    # Split the output back to learnable queries and noisy queries
    Q_N = cat(out[:-N].unsqueeze(0).split(B,dim=2),dim=0) # (G,CxK,B,D)
    Q_L = cat(out[-N:].split(B,dim=1),dim=0) # (GxN,B,D)
    
else:   
    # In inference phase Q_N = none, M = none
    Q = Q_L[:N]
    
out = SA(Q)
\end{lstlisting}
\end{algorithm}


\subsection{Variational Query DeNoising}
\textbf{Challenges with Conventional DeNoising.} While the integration of conventional denoising techniques can stabilize the integration of 3D-attribute into the Bipartite Matching process and offers an initial improvement in model detection performance, this approach encounters the gradient vanishing problem. We can visually demonstrate this issue by inspecting the attention map \(A_i \in \mathbb{R}^{S \times S}\) derived from the Mask Separated Self-Attention mechanism for the \(i\)-th group of learnable queries (\Cref{fig:motivation_vdn}a). Observation of the attention map reveals a significant decoupling effect. Specifically, the attention scores between the noisy queries and the learnable queries approach zero, visibly in the upper-right quadrant of the map. This decoupling is detrimental because it obstructs the efficient backpropagation of gradients originating from the reconstruction loss to the learnable queries, thereby limiting their capacity for further adaptation and improvement.

\begin{figure*}[h!]
    \centering
    \includesvg[width=\textwidth]{images/mono3dv/motivation_vdn.svg}
    \caption{\textbf{The influence of the denoised query on learnable queries.} The self-attention maps trend resulting from the conventional denoising method and the proposed Variational Query DeNoising approach are illustrated in (a) and (b), respectively. (c) presents the average entropy of the attention maps for both methods throughout the entire training period.}
    \label{fig:motivation_vdn}
\end{figure*}

We rigorously examine the underlying causes of the gradient vanishing problem as it pertains to the reconstruction loss's impact on the learnable queries.
As described in \Cref{eq:q_i}, the new set of queries $Q$ is constructed from the sets of learnable queries $Q_L$ and noisy queries $Q_N$.
This newly constructed set $Q$ is then fed into the Mask Separated Self-Attention mechanism. To clearly distinguish the outputs of this mechanism from their inputs, we denote the output by prefacing the input notation with an $o$ at the base letter. The detailed calculations for the Mask Separated Self-Attention are as follows:
\begin{equation}
    \begin{aligned}
        o_{q_i} &= \texttt{concat}\left(o_{q_{N_{i1}}},o_{q_{N_{i2}}},...,o_{q_{N_{iC}}},o_{q_{L_i}}\right) \\
        o_{N_{ijk}} &= \underbrace{\sum_{m=1}^{K} A\left(u,v\right)q_{N_{ijm}}}_{\texttt{noisy queries aggregation}} \\ 
        &+ \underbrace{\sum_{m=1}^{N}A\left(u, w\right)q_{L_{im}}}_{\texttt{learnable queries aggregation}} \\ & \hspace{2em} \texttt{for } j\in \{1,2,...,C\},k\in \{1,2,...,K\} \\
        o_{q_{L_{im}}} &= \sum_{n=1}^N A\left(w, x\right)q_{L_{in}} \\
        & \hspace{3em} \texttt{for } m \in \{1,2,...,N\} \\
        & \texttt{with } 
        u = (j-1)\times K + k,\\
        & \hspace{3.2em} v = (j-1)\times K + m,\\
        & \hspace{3.2em} w = (C-1)\times K + m,\\
        & \hspace{3.2em} x = (C-1)\times K + n\\
    \end{aligned}
    \label{eq:forward_msa}
\end{equation}
where $A \in [0,1]^{S \times S}$ denotes the attention map, $A(u,v)$ refers to the attention score at row $u$ and column $v$.
These outputs are subsequently fed into the remaining part of the decoder and the detection head to compute the loss function. It is important to note that only the outputs of the noisy queries $o_{N_{ijk}}$, are utilized for calculating the reconstruction loss.

\textbf{Backward.} Based on the preceding discussion, the gradient of the reconstruction loss $L_{res}$ with respect to the learnable queries $Q_L$ can be expressed as:
\begin{equation}
    \begin{aligned}
        \frac{\partial L_{res}}{\partial q_{L_{im}}} &= \sum_{j=1}^C\sum_{k=1}^K\frac{\partial L_{res}}{\partial o_{N_{ijk}}} \frac{\partial o_{N_{ijk}}}{\partial q_{L_{im}}} \\
        &= \sum_{j=1}^C\sum_{k=1}^K\frac{\partial L_{res}}{\partial o_{N_{ijk}}} A\left(u,w\right) \\
        & \hspace{1em} \texttt{for } i\in \{1,2,...,G\}, m\in \{1,2,...,N\}
    \end{aligned}
    \label{eq:backward_msa}
\end{equation}
As illustrated in Figure 3(a), in conventional denoising approaches, the attention scores between noisy queries and learnable queries, specifically $A\left(u,w\right)$ , tend to approach zero. This consequently leads to:
\begin{equation}
    \frac{\partial L_{res}}{\partial q_{L_{im}}}\approx 0
\end{equation}
which causes a gradient vanishing problem from the reconstruction loss to the learnable queries.

To quantitatively analyze this behavior, we monitored the sparsity of the attention maps \(A_i\) during the evaluation of Mono3DV across training epochs. Treating attention maps as probability distributions, we employed negative entropy as an intuitive metric for sparsity (a lower entropy indicates greater sparsity). Increased sparsity signifies that noisy queries are predominantly attending to themselves rather than the learnable queries. This elevated self-attention among noisy queries directly diminishes the attention weight (and consequently, the gradient flow) directed towards the learnable queries. As illustrated in \Cref{fig:motivation_vdn}c, the attention map associated with the conventional denoising method (AE) exhibits a rapid decrease in negative entropy (increase in sparsity) throughout training, providing further evidence of the restricted gradient flow and resulting performance plateau.

\noindent\textbf{Variational Query DeNoising.} To address this vanishing gradient issue, a straightforward solution is to increase the attention score $A\left(u,w\right)$. To this end, we propose leveraging the inherent stochastic properties of a Variational Autoencoder (VAE). Our approach aims to introduce beneficial variation into the noisy queries, which in turn increases the entropy of the attention distribution and promotes more robust gradient propagation.

As depicted in the overall architecture shown in \Cref{fig:mono3dv_overall} (Variational Query Generator), the process begins by feeding the initial noisy input boxes through a dedicated box embedding layer. This layer functions as the encoder of our VAE, predicting $\mu$ and $\Sigma$ parameters of a latent distribution. The stochastic noisy queries are then synthesized by sampling from this learned distribution using the reparameterization trick: \(z \sim \mathcal{N}\left(\mu,\Sigma\right)\). These generated stochastic queries are subsequently processed by the model's decoder, analogous to how queries are handled in a standard denoising setup. Training for this denoising process is guided by a denoising loss function, defined as:
\begin{equation}
    L_{DN} = L_{res} + \beta L_{KL}\left(\mathcal{N}\left(\mu,\Sigma\right),\mathcal{N}\left(0,I\right)\right)
    \label{eq:vdn_loss}
\end{equation}
with $L_{res}$ representing reconstruction loss computed from the noisy queries, $L_{KL}$ denotes the Kullback-Leibler divergence loss that regularizes the learned distribution towards a standard normal prior $\mathcal{N}(0,I)$, and $\beta$ is a weighting factor.

\Cref{fig:motivation_vdn}c demonstrates that the proposed Variational Query Denoising (VAE) maintains significantly higher attention map entropy compared to the conventional denoising method, indicating reduced sparsity. Furthermore, as shown in \Cref{fig:motivation_vdn}b (contrast with \Cref{fig:motivation_vdn}a), the attention scores in the upper-right corner of the self-attention map do not converge to zero. This lack of decoupling confirms that the learnable queries continue to effectively interact with and benefit from the denoising process enabled by our variational approach.

\subsection{Loss Function}
\textbf{Forward-Looking Distillation.} We propose that integrating self-distillation enhances the iterative refinement strategy in DETR~\cite{DETR} by transferring knowledge from the final high-performing decoder layer to shallower layers. To prioritize high-quality predictions, we weight the distillation loss using the $\mathbf{IoU}_{3D}$ between the last decoder prediction and the ground-truth 3D box. Following~\cite{SD} success, a shared MLP $\left(f_Q\left(\cdot\right)\right)$ refines the query to improve the quality of the distilled knowledge. The self-distillation loss is formally expressed as:
\begin{equation}
    L_{dis} = \sum_{i=1}^{D-1}{\mathbf{IoU}_{3D}}_{Q_D}\cdot \texttt{Smooth\_L1}\left(f_Q\left(Q_i\right), Q_D\right)
    \label{eq:loss_dis}
\end{equation}
where $Q_i$ represents the output query of the i-th decoder layer, $D$ denotes the total number of decoder layers in the model. It is important to note that Forward-Looking Distillation is applied to both the learnable queries $Q_L$ and the noisy queries $Q_N$. Unlike in the case of learnable queries, the requirement for Hungarian matching between the final decoder layer and ground truth to identify positive queries is obviated for noisy queries. For noisy queries, we instead implement a distillation loss for the queries at each layer, which are derived from the same ground truth. We uniformly applied the reconstruction loss to all bounding box properties, consistent with the formulation of $L_{det}$.

\noindent\textbf{Overall Loss.} The training loss of Mono3DV is composed of three distinct terms: detection loss $L_{det}$, denoising loss $L_{DN}$, and self-distillation loss $L_{dis}$. We adopt $L_{det}$ formulation from MonoDETR~\cite{MonoDETR} that includes losses for object category, projected center point, 2D bounding box, orientation, 3D size, central depth, and depth map. In addition, we propose the denoising loss, $L_{DN}$ derived from Variational Query DeNoising as presented in \Cref{eq:vdn_loss}, and $L_{dis}$ is integrated for the proposed Forward-Looking Distillation as mentioned in \Cref{eq:loss_dis}. The overall loss of Mono3DV is then formulated as:
\begin{equation}
    L_{overall} = \lambda_1L_{det}+ \lambda_{2}L_{DN}+ \lambda_{3}L_{dis}
    \label{eq:overall_loss}
\end{equation}

\subsection{Plug-and-play for Multi-view Detectors}
Besides monocular images, our Variational Query Denoising and Forward-Looking Distillation can also serve as a plug-and-play module upon multi-view methods. Following the baseline MonoDETR~\cite{MonoDETR}, we conduct a multi-view plug-and-play study. This involves integrating Variational Query DeNoising and Forward-Looking Distillation onto two DETR-based multi-view networks: RayDN~\cite{RayDN} and OPEN~\cite{OPEN}. Since these DETR-based multi-view baselines already incorporate a denoising technique, our modification was simplified to adapting the query embedding as a Variational Autoencoder (VAE) and adding a self-distillation loss.